\begin{helper}[Intersecting component shortens the no-repetition case]
Let \((\varphi,T)\) be a minimum bad pair with
\(2<q:=\rho(T)<n:=|E(T)|\).  Suppose distinct colors
\(\varepsilon,\beta\) are missing at \(y_m,y_i\), respectively, with
\(m,i\le q-2\), and \(\beta\) is also missing at \(y_n\).  If the
\((\varepsilon,\beta)\)-component rooted at \(y_n\) meets a displayed
vertex of \(T_{q-1}\), then a smaller bad pair exists, a contradiction.
\end{helper}
\begin{proof}
Let \(Q\) be the stated component.  First suppose that \(Q\) meets the
smaller exact prefix \(T_{q-2}\).  Both colors belong to
\noderef{TreeMissingColors} of that prefix, witnessed by \(y_m\) and
\(y_i\).  The final vertex is outside the prefix by injectivity of the
displayed vertices.  Thus \noderef{TwoColorComponentTashkinovSplice},
applied with \(S=T_{q-2}\) from
\noderef{TreeSequencePrefix}, returns a longer tree sequence \(S'\)
which has \(S\) as an exact initial segment, ends at \(y_n\), and whose
appended exact edges and vertices follow the oriented alternating
component segment.  It also returns a
\noderef{TashkinovAugmentationSequence}; hence
\noderef{TashkinovAugmentationSequencePreservesTashkinovTree}, together
with the Tashkinov clause of \noderef{TashTreeBadPair}, proves that
\(S'\) is a Tashkinov tree under \(\varphi\).

The retained vertices \(y_i,y_m\) are present in \(S\), and the splice
ends at \(y_n\).  In particular, \(\beta\) is missing at the two
distinct vertices \(y_i,y_n\) of \(S'\).  The unchanged \(k\)-triple
and this failure of elementariness show, by
\noderef{TashTreeBadPair}, that \((\varphi,S')\) is bad.  The generic
splice assertion makes consecutive endpoint equations automatic only
after its first appended edge.  Since \(S\) has length \(q-2\), those
equations make \(q-1=(q-2)+1\) a suffix-path start for \(S'\).  Hence
its path index is at most \(q-1<q=\rho(T)\), which is the strict rank
comparison \noderef{TashTreeRankLess}.  This contradicts the universal
minimality clause of \noderef{TashTreeMinimumBadPair}.

It remains to suppose that \(Q\) avoids every displayed vertex of
\(T_{q-2}\) but meets \(T_{q-1}\).  Apply
\noderef{TwoColorComponentTashkinovSplice} with
\(S=T_{q-1}\).  The attachment vertex returned by that theorem lies in
\(Q\cap V(T_{q-1})\).  Exact-prefix vertex identities from
\noderef{TreeSequencePrefix} say that the only vertex of
\(T_{q-1}\) not already in \(T_{q-2}\) is \(y_{q-1}\); therefore the
returned attachment is exactly \(y_{q-1}\).  The special first-edge
equation in the splice statement now says that the first appended edge
joins the last retained vertex \(y_{q-1}\) to the first new vertex.
Together with the later consecutive equations, the entire appended
segment is in path order beginning at position \(q-1\), rather than
only at the generic position \(q\).

As in the first subcase, the retained prefix contains both early
missing-color witnesses, the splice ends at \(y_n\), its augmentation
sequence preserves the Tashkinov property, and the two
missing-\(\beta\) vertices make it bad.  Its path index is at most
\(q-1<q\), so \noderef{TashTreeRankLess} again contradicts the
minimality clause of \noderef{TashTreeMinimumBadPair}.  The two prefix
subcases exhaust the asserted meeting with \(T_{q-1}\).
\end{proof}
